\section{Exercises}\label{chap:heat:ex}

\paragraph{On Section~\ref{ch:heat}.}
\begin{exercises}
\item Show that $c_0c_2-c_1^2=\frac12\sum_{j,l}(a_j-a_l)^2e^{-(a_j+a_l)t}$ and deduce that log convexity of $W$ holds under RH.
\item Recompute the first row of Table~\ref{tab:ch12:hankel} with the first $100$ zeros and explain why no printed digit changes.
\item Show that for a real exponent $a$ the summand of the Williamson measure $U_k$ belonging to $a$ is the inverse gamma density of $ka/\Gamma_k$ with $\Gamma_k\sim$ Gamma$(k,1)$, and compute its variance for $k>2$.
\item\label{ex:ch12:toy} Let $\widetilde W(t)=e^{-t}+2e^{-10t}\cos t$, the heat trace of the three exponents $1$ and $10\pm i$. Show that $\widetilde W>0$ and $\widetilde W'<0$ on $(0,\infty)$, using $\cos t>0$ and $\sin t\ge0$ for $t\le1$ and $2\sqrt{101}\,e^{-9}<1$ for $t\ge1$, so that the first two rungs of Theorem~\ref{thm:ch12:idsd} hold. With $\theta=\arctan\frac1{10}$, show that $\frac{10}{101}<\theta<\frac1{10}$, hence $\frac{2\pi}3<24\theta<\pi$, and deduce that $(-1)^{24}\widetilde W^{(24)}(0)=1+2\cdot101^{12}\cos24\theta<0$, so that $\widetilde W$ is not completely monotone. Then check numerically that the first order $k$ at which $(-1)^k\widetilde W^{(k)}(t)<0$ for some $t>0$ is $12$, near $t=2.87$, and compare with the first failures of Table~\ref{tab:ch13:first}.
\item Show that the term of order $1/u$ of $\varepsilon_\rho$ in Proposition~\ref{prop:ch12:ray} is $\frac{2t}{3u}(\tilde\gamma-u)^3$. For $u=2$, $t=1$ and $k=4$, show that no zero lies in the disc $|\tilde\gamma-u|\le u/2$, so that $c_4(1)=e^{f_0}R$ is carried entirely by the remainder, and compare the bound on $R$ with the contribution of the first zero.
\item Show that $m(0)>-6$ and, comparing the series for $\psi$ with an integral, that $m(60)>1$. Deduce that for $t\le10^{-6}$ the archimedean term of \eqref{eq:ch2:Warith} is at least $1/(4\sqrt{\pi t})-210/\pi>55$ and the prime term at most $(\pi t)^{-1/2}e^{-(\log2)^2/4t}<0.01$. With the argument of the proof of Theorem~\ref{thm:ch12:idsd} for $t\ge\tau$, taken at height $10^4$ and from $t=10^{-6}$, conclude that $W>0$ on $(0,\infty)$ follows from the verification of RH below height $10^4$ by estimates that can all be checked by hand.
\end{exercises}

\paragraph{On Section~\ref{ch:arith}.}
\begin{exercises}
\item Verify the arithmetic form of $W(t)$ at $t=0.01$ with the prime powers up to $10^4$, and explain why so few suffice by showing that only $n=2$ and $n=3$ affect the first fifteen digits.
\item\label{ex:ch14:law} Show that $m$ increases on $[0,\infty)$ and vanishes at a single point $r_*$, and deduce that $e^{r_*^2t}A(t)$ decreases in $t$, so that $A$ changes sign at most once; with Table~\ref{tab:ch14:W}, conclude that $A<0$ for $t\ge0.005$. Then show that for $t\ge0.05$ the number of digits lost in the table equals $\log_{10}2+(\gamma_1^2+\frac14)t/\log10$ up to an exponentially small error, and identify the source of each term.
\item Show that the regularised constants of the gamma, pole and prime parts of the variogram form sum to zero if and only if $\xi'(\frac12)=0$.
\end{exercises}

\paragraph{On Section~\ref{ch:spacing}.}
\begin{exercises}
\item\label{ex:ch15:odd} Let $\ell=\Xi'/\Xi$. Show that $\xi(\frac12+\varepsilon+ib)=\Xi(b-i\varepsilon)$, and deduce from the Taylor series of $\log\Xi$ about $b$ that $\mu_\varepsilon(C_k)-1=-\frac\varepsilon\pi[\ell(b_k)-\ell(b_{k-1})]+\frac{\varepsilon^3}{6\pi}[\ell''(b_k)-\ell''(b_{k-1})]+O(\varepsilon^5)$. Under RH, check that the term in $\varepsilon^3$ agrees with the expansion of $F_\varepsilon$ to third order.
\item Prove the arctangent formulas for the cell masses of the sinh benchmark, and deduce that its cell distance for $N$ atoms is $(4-1/(N+\frac12))\,\varepsilon/\pi+O(\varepsilon^3)$, and $(\frac43-1/(N+\frac12))\,\varepsilon/\pi+O(\varepsilon^3)$ when the first cell is $(0,\frac32]$.
\item Show that the Wigner surmise $p_{\rm W}$ has mean $1$ and variance $3\pi/8-1$, and that $\E s^{-2}=8/\pi$ under it, while $\E s^{-2}=\infty$ for the Poisson law.
\item Show that $\sum_{m\ne0}(m\Delta)^{-2}=\pi^2/(3\Delta^2)$, deduce the lattice threshold $(\sqrt6/\pi)\Delta$ of the visibility condition, and find the height above which the threshold is less than $\frac12$.
\item\label{ex:ch15:beside} In a local model, put an off line pair at distance $d=0.3$ from the line at the ordinate $\gamma$ and zeros on the line at $\gamma+0.2+m$ and $\gamma-4-m$ for $m\ge0$, and ignore the reflected kernels. At $\varepsilon=0.005$, show that $2/d^2<\sum_j\delta_j^{-2}$, so that the visibility condition predicts a positive density at $\gamma$, and check that $\rho_\varepsilon(\gamma)>0$ while $\rho_\varepsilon(\gamma-0.079)<0$.
\end{exercises}
